%% socket.tex
%%
%% Copyright (C) 2000-2018 Simone Piccardi.  Permission is granted to
%% copy, distribute and/or modify this document under the terms of the GNU Free
%% Documentation License, Version 1.1 or any later version published by the
%% Free Software Foundation; with the Invariant Sections being "Un preambolo",
%% with no Front-Cover Texts, and with no Back-Cover Texts.  A copy of the
%% license is included in the section entitled "GNU Free Documentation
%% License".
%%

\chapter{I socket}
\label{cha:socket_intro}

In questo capitolo inizieremo a spiegare le caratteristiche salienti della
principale interfaccia per la programmazione di rete, quella dei
\textit{socket}, che, pur essendo nata in ambiente Unix, è usata ormai da
tutti i sistemi operativi.

Dopo una breve panoramica sulle caratteristiche di questa interfaccia vedremo
come creare un socket e come collegarlo allo specifico protocollo di rete che
si utilizzerà per la comunicazione. Per evitare un'introduzione puramente
teorica concluderemo il capitolo con un primo esempio di applicazione.

\section{Introduzione ai socket}
\label{sec:sock_overview}

In questa sezione daremo descrizione essenziale di cosa sono i \textit{socket}
e di quali sono i concetti fondamentali da tenere presente quando si ha a che
fare con essi; ne illustreremo poi le caratteristiche e le differenti
tipologie presenti ed infine tratteremo le modalità con cui possono essere
creati.

\index{socket!definizione|(}

\subsection{Cosa sono i \textit{socket}}
\label{sec:sock_socket_def}

I \textit{socket} (una traduzione letterale potrebbe essere \textsl{presa}, ma
essendo universalmente noti come \textit{socket} utilizzeremo sempre la parola
inglese) sono uno dei principali meccanismi di comunicazione utilizzato in
ambito Unix, e li abbiamo brevemente incontrati in
sez.~\ref{sec:ipc_socketpair}, fra i vari meccanismi di intercomunicazione fra
processi. 

Un socket costituisce in sostanza un canale di comunicazione fra due processi
su cui si possono leggere e scrivere dati analogo a quello di una
\textit{pipe} (vedi sez.~\ref{sec:ipc_pipes}) ma, a differenza di questa e
degli altri meccanismi esaminati nel capitolo cap.~\ref{cha:IPC}, i socket non
sono limitati alla comunicazione fra processi che girano sulla stessa
macchina, ma possono realizzare la comunicazione anche attraverso la rete.

Quella dei socket costituisce infatti la principale interfaccia usata nella
programmazione di rete.  La loro origine risale al 1983, quando furono
introdotti in BSD 4.2; l'interfaccia è rimasta sostanzialmente la stessa, con
piccole modifiche, negli anni successivi. Benché siano state sviluppate
interfacce alternative, originate dai sistemi SVr4 come la XTI (\textit{X/Open
  Transport Interface}) nessuna ha mai raggiunto la diffusione e la popolarità
di quella dei socket (né tantomeno la stessa usabilità e flessibilità) ed oggi
sono praticamente dimenticate.

La flessibilità e la genericità dell'interfaccia inoltre consente di
utilizzare i socket con i più disparati meccanismi di comunicazione, e non
solo con l'insieme dei protocolli TCP/IP, anche se questa sarà comunque quella
di cui tratteremo in maniera più estesa.

Per capire il funzionamento dei socket occorre avere presente il funzionamento
dei protocolli di rete che su utilizzeranno (ed in particolare quelli del
TCP/IP già illustrati in sez.~\ref{sec:net_tpcip}), ma l'interfaccia è del
tutto generale e benché le problematiche, e quindi le modalità di risolvere i
problemi, siano diverse a seconda del tipo di protocollo di comunicazione
usato, le funzioni da usare nella gestione dei socket restano le stesse.

Per questo motivo una semplice descrizione dell'interfaccia è assolutamente
inutile, in quanto il comportamento di quest'ultima e le problematiche da
affrontare cambiano radicalmente a seconda del tipo di comunicazione usato.
La scelta di questo tipo di comunicazione (sovente anche detto \textsl{stile})
va infatti ad incidere sulla semantica che verrà utilizzata a livello utente
per gestire la comunicazione cioè su come inviare e ricevere i dati e sul
comportamento effettivo delle funzioni utilizzate.

La scelta di uno \textsl{stile} dipende sia dai meccanismi disponibili, sia
dal tipo di comunicazione che si vuole effettuare. Ad esempio alcuni tipi di
comunicazione considerano i dati come una sequenza continua di byte, in quello
che viene chiamato un \textsl{flusso} (in inglese \textit{stream}), mentre
altri invece li raggruppano in \textsl{pacchetti} (in inglese
\textit{datagram}) che vengono sempre inviati in blocchi separati e non
divisibili.

Un altro esempio delle differenze fra i diversi tipi di comunicazione concerne
la possibilità che essa possa o meno perdere dati nella trasmissione, che
possa o meno rispettare l'ordine in cui i dati inviati e ricevuti, o che possa
accadere di inviare dei pacchetti di dati più volte (differenze che ad esempio
sono presenti nel caso di utilizzo dei protocolli TCP o UDP).

Un terzo esempio di differenza nel tipo di comunicazione concerne il modo in
cui essa avviene nei confronti dei corrispondenti, in certi casi essa può
essere condotta con una connessione diretta con un solo corrispondente, come
per una telefonata; altri casi possono prevedere una comunicazione come per
lettera, in cui si scrive l'indirizzo su ogni pacchetto, altri ancora una
comunicazione uno a molti come il \textit{broadcast} ed il \textit{multicast},
in cui i pacchetti possono venire emessi su appositi ``\textsl{canali}'' dove
chiunque si collega possa riceverli.

É chiaro che ciascuno di questi diversi aspetti è associato ad un tipo di
comunicazione che comporta una modalità diversa di gestire la stessa, ad
esempio se la comunicazione è inaffidabile occorrerà essere in grado di
gestire la perdita o il rimescolamento dei dati, se è a pacchetti questi
dovranno essere opportunamente trattati, se è uno a molti occorrerà tener
conto della eventuale unidirezionalità della stessa, ecc.

\index{socket!definizione|)}


\subsection{La creazione di un socket}
\label{sec:sock_creation}

Come accennato l'interfaccia dei socket è estremamente flessibile e permette
di interagire con protocolli di comunicazione anche molto diversi fra di loro;
in questa sezione vedremo come è possibile creare un socket e come specificare
il tipo di comunicazione che esso deve utilizzare.

La creazione di un socket avviene attraverso l'uso della funzione di sistema 
\funcd{socket}; essa restituisce un \textit{file descriptor} (del tutto
analogo a quelli che si ottengono per i file di dati e le \textit{pipe},
descritti in sez.~\ref{sec:file_fd}) che serve come riferimento al socket; il
suo prototipo è:

\begin{funcproto}{
\fhead{sys/socket.h}
\fdecl{int socket(int domain, int type, int protocol)}
\fdesc{Apre un socket.} 
}

{La funzione ritorna un valore positivo in caso di successo e $-1$ per un
  errore, nel qual caso \var{errno} assumerà uno dei valori:
  \begin{errlist}
  \item[\errcode{EACCES}] non si hanno privilegi per creare un socket nel
    dominio o con il protocollo specificato.
  \item[\errcode{EAFNOSUPPORT}] famiglia di indirizzi non supportata.
  \item[\errcode{EINVAL}] argomento \param{type} invalido.
  \item[\errcode{EMFILE}] si è ecceduta la tabella dei file.
  \item[\errcode{ENFILE}] si è raggiunto il limite massimo di file aperti.
  \item[\errcode{ENOBUFS}] non c'è sufficiente memoria per creare il socket
    (può essere anche \errval{ENOMEM}).
  \item[\errcode{EPROTONOSUPPORT}] il tipo di socket o il protocollo scelto
    non sono supportati nel dominio.
  \end{errlist}
  ed inoltre a seconda del protocollo usato, potranno essere generati altri
  errori, che sono riportati nelle pagine di manuale relative al protocollo.}
\end{funcproto}


La funzione ha tre argomenti, \param{domain} specifica il dominio del socket
(definisce cioè, come vedremo in sez.~\ref{sec:sock_domain}, la famiglia di
protocolli usata), \param{type} specifica il tipo di socket (definisce cioè,
come vedremo in sez.~\ref{sec:sock_type}, lo stile di comunicazione) e
\param{protocol} il protocollo; in genere quest'ultimo è indicato
implicitamente dal tipo di socket, per cui di norma questo valore viene messo
a zero (con l'eccezione dei \textit{raw socket}).

% TODO: l'ultimo argomento viene usato anche dai nuovi ping socket introdotti
% con il kernel 3.0, vedi anche http://lwn.net/Articles/420799/ e
% http://git.kernel.org/?p=linux/kernel/git/torvalds/linux-2.6.git;a=commitdiff;h=c319b4d76b9e583a5d88d6bf190e079c4e43213d 

Si noti che la creazione del socket si limita ad allocare le opportune
strutture nel kernel (sostanzialmente una voce nella \textit{file table}) e
non comporta nulla riguardo all'indicazione degli indirizzi remoti o locali
attraverso i quali si vuole effettuare la comunicazione. Questo significa che
la funzione da sola non è in grado di fornire alcun tipo di comunicazione. 


\subsection{Il dominio dei socket}
\label{sec:sock_domain}

Dati i tanti e diversi protocolli di comunicazione disponibili, esistono vari
tipi di socket, che vengono classificati raggruppandoli in quelli che si
chiamano \textsl{domini}.  La scelta di un dominio equivale in sostanza alla
scelta di una famiglia di protocolli, e viene effettuata attraverso
l'argomento \param{domain} della funzione \func{socket}. Ciascun dominio ha un
suo nome simbolico che convenzionalmente è indicato da una costante che inizia
per \texttt{PF\_}, sigla che sta per \textit{protocol family}, altro nome con
cui si indicano i domini.

A ciascun tipo di dominio corrisponde un analogo nome simbolico, anch'esso
associato ad una costante, che inizia invece per \texttt{AF\_} (da
\textit{address family}) che identifica il formato degli indirizzi usati in
quel dominio. Le pagine di manuale di Linux si riferiscono a questi indirizzi
anche come \textit{name space}, (nome che invece il manuale della \acr{glibc}
riserva a quello che noi abbiamo chiamato domini) dato che identificano il
formato degli indirizzi usati in quel dominio per identificare i capi della
comunicazione.

\begin{table}[!htb]
  \footnotesize
  \centering
  \begin{tabular}[c]{|l|l|l|l|}
       \hline
       \textbf{Nome}&\textbf{Valore}&\textbf{Utilizzo}&\textbf{Man page} \\
       \hline
       \hline
       \constd{AF\_UNSPEC}   & 0& Non specificato               &            \\
       \constd{AF\_LOCAL}    & 1& Local communication           & unix(7)    \\
       \constd{AF\_UNIX}, \constd{AF\_FILE}&1&Sinonimi di \const{AF\_LOCAL}& \\
       \constd{AF\_INET}     & 2& IPv4 Internet protocols       & ip(7)      \\
       \constd{AF\_AX25}     & 3& Amateur radio AX.25 protocol  &            \\
       \constd{AF\_IPX}      & 4& IPX - Novell protocols        &            \\
       \constd{AF\_APPLETALK}& 5& Appletalk                     & ddp(7)     \\
       \constd{AF\_NETROM}   & 6& Amateur radio NetROM          &            \\
       \constd{AF\_BRIDGE}   & 7& Multiprotocol bridge          &            \\
       \constd{AF\_ATMPVC}   & 8& Access to raw ATM PVCs        &            \\
       \constd{AF\_X25}      & 9& ITU-T X.25 / ISO-8208 protocol& x25(7)     \\
       \constd{AF\_INET6}    &10& IPv6 Internet protocols       & ipv6(7)    \\
       \constd{AF\_ROSE}     &11& Amateur Radio X.25 PLP        &            \\
       \constd{AF\_DECnet}   &12& Reserved for DECnet project   &            \\
       \constd{AF\_NETBEUI}  &13& Reserved for 802.2LLC project &            \\
       \constd{AF\_SECURITY} &14& Security callback pseudo AF   &            \\
       \constd{AF\_KEY}      &15& AF\_KEY key management API    &            \\
       \constd{AF\_NETLINK}  &16& Kernel user interface device  & netlink(7) \\
       \constd{AF\_ROUTE}    &16& Sinonimo di \const{AF\_NETLINK} emula BSD.&\\
       \constd{AF\_PACKET}   &17& Low level packet interface    & packet(7)  \\
       \constd{AF\_ASH}      &18& Ash                           &    \\
       \constd{AF\_ECONET}   &19& Acorn Econet                  &    \\
       \constd{AF\_ATMSVC}   &20& ATM SVCs                      &    \\
       \constd{AF\_RDS}      &21& RDS Sockets                   &    \\
       \constd{AF\_SNA}      &22& Linux SNA Project             &    \\
       \constd{AF\_IRDA}     &23& IRDA socket (infrarossi)      & irda(7)    \\
       \constd{AF\_PPPOX}    &24& PPPoX socket                  &    \\
       \constd{AF\_WANPIPE}  &25& Wanpipe API socket            &    \\
       \constd{AF\_LLC}      &26& Linux LLC                     &    \\
       \constd{AF\_IB}       &27&  Native InfiniBand address    &    \\
       \constd{AF\_MPLS}     &28& MPSL                          &    \\
       \constd{AF\_CAN}      &29& Controller Are Network        &    \\
       \constd{AF\_TIPC}     &30& TIPC sockets                  &    \\
       \constd{AF\_BLUETOOTH}&31& Bluetooth socket              &    \\
       \constd{AF\_IUCV}     &32& IUCV sockets                  &    \\
       \constd{AF\_RXRPC}    &33& RxRPC sockets                 &    \\
       \constd{AF\_ISDN}     &34& mISDN sockets                 &    \\
       \constd{AF\_PHONET}   &35& Phonet sockets                &    \\
       \constd{AF\_IEEE802154}&36& IEEE802154 sockets           &    \\
       \constd{AF\_CAIF}     &37& CAIF sockets                  &    \\
       \constd{AF\_ALG}      &38& Algorithm sockets             &    \\
       \constd{AF\_NFC}      &39& NFC sockets                   &    \\
       \constd{AF\_VSOCK}    &40& vSockets                      &    \\
       \constd{AF\_KCM}      &41& Kernel Connection Multiplexor &    \\
       \constd{AF\_QIPCRTR}  &42& Qualcomm IPC Router           &    \\
       \constd{AF\_SMC}      &43& smc sockets                   &    \\
       \hline
  \end{tabular}
  \caption{Famiglie di protocolli definiti in Linux.} 
  \label{tab:net_pf_names}
\end{table}

L'idea alla base della distinzione fra questi due insiemi di costanti era che
una famiglia di protocolli potesse supportare vari tipi di indirizzi, per cui
il prefisso \texttt{PF\_} si sarebbe dovuto usare nella creazione dei socket e
il prefisso \texttt{AF\_} in quello delle strutture degli indirizzi. Questo è
quanto specificato anche dallo standard POSIX.1g, ma non esistono a tuttora
famiglie di protocolli che supportino diverse strutture di indirizzi, per cui
nella pratica questi due nomi sono equivalenti e corrispondono agli stessi
valori numerici.\footnote{in Linux, come si può verificare andando a guardare
  il contenuto di \file{bits/socket.h}, le costanti sono esattamente le stesse
  e ciascuna \texttt{AF\_} è definita alla corrispondente \texttt{PF\_} e con
  lo stesso nome.} Qui si sono indicati i nomi con il prefisso \texttt{AF\_}
seguendo la convenzione usata nelle pagine di manuale.

I domini (e i relativi nomi simbolici), così come i nomi delle famiglie di
indirizzi, sono definiti dall'\textit{header file} \headfiled{socket.h}. Un
elenco, aggiornato alla versione 4.15, delle famiglie di protocolli
disponibili in Linux è riportato in tab.~\ref{tab:net_pf_names}.  L'elenco
indica tutti i protocolli definiti; fra questi però saranno utilizzabili solo
quelli per i quali si è compilato il supporto nel kernel (o si sono caricati
gli opportuni moduli), viene definita anche una costante \constd{AF\_MAX} che
indica il valore massimo associabile ad un dominio.

Si tenga presente che non tutte le famiglie di protocolli sono utilizzabili
dall'utente generico, ad esempio in generale tutti i socket di tipo
\const{SOCK\_RAW} possono essere creati solo da processi che hanno i privilegi
di amministratore (cioè con \ids{UID} effettivo uguale a zero) o dotati della
\textit{capability} \const{CAP\_NET\_RAW}.


\subsection{Il tipo di socket}
\label{sec:sock_type}

La scelta di un dominio non comporta però la scelta dello stile di
comunicazione, questo infatti viene a dipendere dal protocollo che si andrà ad
utilizzare fra quelli disponibili nella famiglia scelta. L'interfaccia dei
socket permette di scegliere lo stile di comunicazione indicando il tipo di
socket con l'argomento \param{type} di \func{socket}. Linux mette a
disposizione vari tipi di socket (che corrispondono a quelli che il manuale
della \acr{glibc} \cite{GlibcMan} chiama \textit{styles}) identificati dalle
seguenti costanti:\footnote{le pagine di manuale POSIX riportano solo i primi
  tre tipi, Linux supporta anche gli altri, come si può verificare nel file
  \texttt{include/linux/net.h} dei sorgenti del kernel.}

\begin{basedescript}{\desclabelwidth{2.0cm}\desclabelstyle{\nextlinelabel}}
\item[\constd{SOCK\_STREAM}] Provvede un canale di trasmissione dati
  bidirezionale, sequenziale e affidabile. Opera su una connessione con un
  altro socket. I dati vengono ricevuti e trasmessi come un flusso continuo di
  byte (da cui il nome \textit{stream}) e possono essere letti in blocchi di
  dimensioni qualunque. Può supportare la trasmissione dei cosiddetti dati
  urgenti (o \textit{out-of-band}, vedi sez.~\ref{sec:TCP_urgent_data}).
\item[\constd{SOCK\_DGRAM}] Viene usato per trasmettere pacchetti di dati
  (\textit{datagram}) di lunghezza massima prefissata, indirizzati
  singolarmente. Non esiste una connessione e la trasmissione è effettuata in
  maniera non affidabile.
\item[\constd{SOCK\_SEQPACKET}] Provvede un canale di trasmissione di dati
  bidirezionale, sequenziale e affidabile. Opera su una connessione con un
  altro socket. I dati possono vengono trasmessi per pacchetti di dimensione
  massima fissata, e devono essere letti integralmente da ciascuna chiamata a
  \func{read}.
\item[\constd{SOCK\_RAW}] Provvede l'accesso a basso livello ai protocolli di
  rete e alle varie interfacce. I normali programmi di comunicazione non
  devono usarlo, è riservato all'uso di sistema.
\item[\constd{SOCK\_RDM}] Provvede un canale di trasmissione di dati
  affidabile, ma in cui non è garantito l'ordine di arrivo dei pacchetti.
\item[\constd{SOCK\_PACKET}] Obsoleto, non deve essere più usato (e pertanto
  non ne parleremo ulteriormente).
\end{basedescript}

A partire dal kernel 2.6.27 l'argomento \param{type} della funzione
\func{socket} assume un significato ulteriore perché può essere utlizzato per
impostare dei flag relativi alle caratteristiche generali del \textit{socket}
non strettamente attinenti all'indicazione del tipo secondo i valori appena
illustrati. Essi infatti possono essere combinati con un OR aritmetico delle
ulteriori costanti:
\begin{basedescript}{\desclabelwidth{2.5cm}\desclabelstyle{\nextlinelabel}}
\item[\constd{SOCK\_CLOEXEC}] imposta il flag di \textit{close-on-exec} sul
  file descriptor del socket, ottenendo lo stesso effetto del flag
  \const{O\_CLOEXEC} di \func{open} (vedi tab.~\ref{tab:open_operation_flag}),
  di cui costituisce l'analogo.

\item[\constd{SOCK\_NONBLOCK}] crea il socket in modalità non-bloccante, con
  effetti identici ad una successiva chiamata a \func{fcntl} per impostare il
  flag di \const{O\_NONBLOCK} sul file descriptor (si faccia di nuovo
  riferimenti al significato di quest'ultimo come spiegato in
  tab.~\ref{tab:open_operation_flag}).
\end{basedescript}

Si tenga presente inoltre che non tutte le combinazioni fra una famiglia di
protocolli e un tipo di socket sono valide, in quanto non è detto che in una
famiglia esista un protocollo per ciascuno dei diversi stili di comunicazione
appena elencati.

\begin{table}[htb]
  \footnotesize
  \centering
  \begin{tabular}{|l|c|c|c|c|c|}
    \hline
    \multicolumn{1}{|c|}{\textbf{Famiglia}}&
    \multicolumn{5}{|c|}{\textbf{Tipo}}\\
    \hline
    \hline
    &\const{SOCK\_STREAM} &\const{SOCK\_DGRAM}     &\const{SOCK\_RAW}& 
      \const{SOCK\_RDM}&\const{SOCK\_SEQPACKET} \\
     \hline
    \const{AF\_UNIX}     &  si & si  &  --  & -- &  si\footnotemark \\
     \hline
    \const{AF\_LOCAL}&\multicolumn{5}{|c|}{sinonimo di \const{AF\_UNIX}}\\
     \hline
    \const{AF\_INET}      & TCP & UDP & IPv4 & --  & --  \\
     \hline
    \const{AF\_INET6}     & TCP & UDP & IPv6 & --  & -- \\
     \hline
    \const{AF\_IPX}       & --  & si  &  --  & --  & -- \\
     \hline
    \const{AF\_NETLINK}   & --  & si  &  si  & --  & -- \\
     \hline
    \const{AF\_X25}       & --  & --  &  --  & --  & si \\
     \hline
    \const{AF\_AX25}      & --  & si  &  si  & --  & si \\
     \hline
%    \const{AF\_ATMPVC}    & --  & --  &  --  & --  & -- \\
%     \hline
    \const{AF\_APPLETALK} & --  & si  &  si  & --  & -- \\
     \hline
    \const{AF\_PACKET}    & --  & si  &  si  & --  & -- \\
     \hline
    \const{AF\_KEY}       & --  & --  &  si  & --  & -- \\
     \hline
    \const{AF\_IRDA}      & si  & si  &  si  & --  & si \\
     \hline
    \const{AF\_NETROM}    & --  & --  &  --  & --  & si \\
     \hline
    \const{AF\_ROSE}      & --  & --  &  --  & --  & si \\
     \hline
    \const{AF\_RDS}       & --  & --  &  --  & --  & si \\
     \hline
    \const{AF\_ECONET}    & --  & si  &  --  & --  & -- \\
     \hline
  \end{tabular}
  \caption{Combinazioni valide di dominio e tipo di protocollo per la 
    funzione \func{socket}.}
  \label{tab:sock_sock_valid_combinations}
\end{table}

\footnotetext{supportati a partire dal kernel 2.6.4 per socket che conservano
  i limiti dei messaggi e li consegnano in sequenza ordinata.}

In tab.~\ref{tab:sock_sock_valid_combinations} sono mostrate le combinazioni
valide possibili per le principali famiglie di protocolli. Per ogni
combinazione valida si è indicato il tipo di protocollo, o la parola
\textsl{si} qualora il protocollo non abbia un nome definito, mentre si sono
lasciate vuote le caselle per le combinazioni non supportate.


\section{Le strutture degli indirizzi dei socket}
\label{sec:sock_sockaddr}

Come si è visto nella creazione di un socket non si specifica nulla oltre al
tipo di famiglia di protocolli che si vuole utilizzare, in particolare nessun
indirizzo che identifichi i due capi della comunicazione. La funzione infatti
si limita ad allocare nel kernel quanto necessario per poter poi realizzare la
comunicazione.

Gli indirizzi infatti vengono specificati attraverso apposite strutture che
vengono utilizzate dalle altre funzioni della interfaccia dei socket, quando
la comunicazione viene effettivamente realizzata.  Ogni famiglia di protocolli
ha ovviamente una sua forma di indirizzamento e in corrispondenza a questa una
sua peculiare struttura degli indirizzi. I nomi di tutte queste strutture
iniziano per \var{sockaddr\_}; quelli propri di ciascuna famiglia vengono
identificati dal suffisso finale, aggiunto al nome precedente.


\subsection{La struttura generica}
\label{sec:sock_sa_gen}

Le strutture degli indirizzi vengono sempre passate alle varie funzioni
attraverso puntatori (cioè \textit{by reference}), ma le funzioni devono poter
maneggiare puntatori a strutture relative a tutti gli indirizzi possibili
nelle varie famiglie di protocolli; questo pone il problema di come passare
questi puntatori, il C moderno risolve questo problema coi i puntatori
generici (i \ctyp{void *}), ma l'interfaccia dei socket è antecedente alla
definizione dello standard ANSI C, e per questo nel 1982 fu scelto di definire
una struttura generica per gli indirizzi dei socket, \struct{sockaddr}, che si
è riportata in fig.~\ref{fig:sock_sa_gen_struct}.

\begin{figure}[!htb]
  \footnotesize \centering
  \begin{minipage}[c]{0.80\textwidth}
    \includestruct{listati/sockaddr.h}
  \end{minipage} 
  \caption{La struttura generica degli indirizzi dei socket
    \structd{sockaddr}.} 
  \label{fig:sock_sa_gen_struct}
\end{figure}

Tutte le funzioni dei socket che usano gli indirizzi sono definite usando nel
prototipo un puntatore a questa struttura; per questo motivo quando si
invocano dette funzioni passando l'indirizzo di un protocollo specifico
occorrerà eseguire una conversione del relativo puntatore.

I tipi di dati che compongono la struttura sono stabiliti dallo standard
POSIX.1g e li abbiamo riassunti in tab.~\ref{tab:sock_data_types} con i
rispettivi file di include in cui sono definiti; la struttura è invece
definita nell'include file \headfile{sys/socket.h}.

\begin{table}[!htb]
  \centering
  \footnotesize
  \begin{tabular}{|l|l|l|}
    \hline
    \multicolumn{1}{|c|}{\textbf{Tipo}}& 
    \multicolumn{1}{|c|}{\textbf{Descrizione}}& 
    \multicolumn{1}{|c|}{\textbf{Header}} \\
    \hline
    \hline
    \typed{int8\_t}   & intero a 8 bit con segno   & \headfile{sys/types.h}\\
    \typed{uint8\_t}  & intero a 8 bit senza segno & \headfile{sys/types.h}\\
    \typed{int16\_t}  & intero a 16 bit con segno  & \headfile{sys/types.h}\\
    \typed{uint16\_t} & intero a 16 bit senza segno& \headfile{sys/types.h}\\
    \typed{int32\_t}  & intero a 32 bit con segno  & \headfile{sys/types.h}\\
    \typed{uint32\_t} & intero a 32 bit senza segno& \headfile{sys/types.h}\\
    \hline
    \typed{sa\_family\_t} & famiglia degli indirizzi&\headfile{sys/socket.h}\\
    \typed{socklen\_t} & lunghezza (\type{uint32\_t}) dell'indirizzo di
    un socket& \headfile{sys/socket.h}\\
    \hline
    \typed{in\_addr\_t} & indirizzo IPv4 (\type{uint32\_t}) & 
    \headfile{netinet/in.h}\\
    \typed{in\_port\_t} & porta TCP o UDP (\type{uint16\_t})& 
    \headfile{netinet/in.h}\\
    \hline
  \end{tabular}
  \caption{Tipi di dati usati nelle strutture degli indirizzi, secondo quanto 
    stabilito dallo standard POSIX.1g.}
  \label{tab:sock_data_types}
\end{table}

In alcuni sistemi la struttura è leggermente diversa e prevede un primo membro
aggiuntivo \code{uint8\_t sin\_len} (come riportato da R. Stevens in
\cite{UNP1}). Questo campo non verrebbe usato direttamente dal programmatore e
non è richiesto dallo standard POSIX.1g, in Linux pertanto non esiste. Il
campo \type{sa\_family\_t} era storicamente un \ctyp{unsigned short}.

Dal punto di vista del programmatore l'unico uso di questa struttura è quello
di fare da riferimento per il casting, per il kernel le cose sono un po'
diverse, in quanto esso usa il puntatore per recuperare il campo
\var{sa\_family}, comune a tutte le famiglie, con cui determinare il tipo di
indirizzo; per questo motivo, anche se l'uso di un puntatore \ctyp{void *}
sarebbe più immediato per l'utente (che non dovrebbe più eseguire il casting),
è stato mantenuto l'uso di questa struttura.

Se si usa una struttura \struct{sockaddr} per allocare delle variabili
generiche da usare in seguito per degli indirizzi si pone il problema che
niente assicura che i dati necessari per le varie famiglie di indirizzi
possano rientrare nella dimensione del campo \var{sa\_data} indicata in
fig.~\ref{fig:sock_sa_gen_struct}, anzi, come vedremo in
sez.~\ref{sec:sock_sa_ipv6}, nel caso di indirizzi IPv6 questa non è proprio
sufficiente. 

\begin{figure}[!htb]
  \footnotesize \centering
  \begin{minipage}[c]{0.95\textwidth}
    \includestruct{listati/sockaddr_storage.h}
  \end{minipage} 
  \caption{La struttura generica degli indirizzi dei socket
    \structd{sockaddr\_storage}.} 
  \label{fig:sock_sa_storage_struct}
\end{figure}

Per questo l'interfaccia di programmazione dei socket prevede la defizione di
una speciale struttura \struct{sockaddr\_storage} illustrata in
fig.~\ref{fig:sock_sa_storage_struct}, in cui di nuovo si usa il primo campo
(\var{ss\_family}) per indicare il tipo di indirizzo, ed in cui i campi
successivi sono utilizzati per allineare i dati al tipo di architettura
hardware utilizzata, e per allocare uno spazio sufficiente ampio per contenere
qualunque tipo di indirizzo supportato. Allocando questa struttura si ha la
certezza di non eccedere le dimensioni qualunque sia il tipo di indirizzi che
si useranno, pertanto risulta utile tutte le volte che si devono gestire in
maniera generica tipi di indirizzi diversi (ad esempio IPv4 ed IPv6).


\subsection{La struttura degli indirizzi IPv4}
\label{sec:sock_sa_ipv4}

I socket di tipo \const{AF\_INET} vengono usati per la comunicazione
attraverso Internet; la struttura per gli indirizzi per un socket Internet (se
si usa IPv4) è definita come \struct{sockaddr\_in} nell'header file
\headfiled{netinet/in.h} ed ha la forma mostrata in
fig.~\ref{fig:sock_sa_ipv4_struct}, conforme allo standard POSIX.1g.

\begin{figure}[!htb]
  \footnotesize\centering
  \begin{minipage}[c]{0.80\textwidth}
    \includestruct{listati/sockaddr_in.h}
  \end{minipage} 
  \caption{La struttura \structd{sockaddr\_in} degli indirizzi dei socket
    Internet (IPv4) e la struttura \structd{in\_addr} degli indirizzi IPv4.}
  \label{fig:sock_sa_ipv4_struct}
\end{figure}

L'indirizzo di un socket Internet (secondo IPv4) comprende l'indirizzo
Internet di un'interfaccia più un \textsl{numero di porta} (affronteremo in
dettaglio il significato di questi numeri in sez.~\ref{sec:TCP_port_num}).  Il
protocollo IP di per sé non prevede numeri di porta, questi sono utilizzati
solo dai protocolli di livello superiore come TCP e UDP, ma devono essere
indicati qui. Inoltre questa struttura viene usata anche per i socket RAW che
accedono direttamente al livello di IP, in questo caso il numero della porta
deve essere impostato al numero di protocollo.

Il membro \var{sin\_family} deve essere sempre impostato a \constd{AF\_INET},
altrimenti si avrà un errore di \errcode{EINVAL}; il membro \var{sin\_port}
specifica il \textsl{numero di porta}. I numeri di porta sotto il 1024 sono
chiamati \textsl{riservati} in quanto utilizzati da servizi standard e
soltanto processi con i privilegi di amministratore (con \ids{UID} effettivo
uguale a zero) o con la \textit{capability} \const{CAP\_NET\_BIND\_SERVICE}
possono usare la funzione \func{bind} (che vedremo in
sez.~\ref{sec:TCP_func_bind}) su queste porte.

Il membro \var{sin\_addr} contiene un indirizzo Internet, e viene acceduto sia
come struttura (un resto di una implementazione precedente in cui questa era
una \dirct{union} usata per accedere alle diverse classi di indirizzi) che
direttamente come intero. In \headfile{netinet/in.h} vengono definite anche
alcune costanti che identificano alcuni indirizzi speciali, riportati in
tab.~\ref{tab:TCP_ipv4_addr}, che rincontreremo più avanti.

Infine occorre sottolineare che sia gli indirizzi che i numeri di porta devono
essere specificati in quello che viene chiamato \textit{network order}, cioè
con i bit ordinati in formato \textit{big endian} (vedi
sez.~\ref{sec:endianness}), questo comporta la necessità di usare apposite
funzioni di conversione per mantenere la portabilità del codice (vedi
sez.~\ref{sec:sock_addr_func} per i dettagli del problema e le relative
soluzioni).


\subsection{La struttura degli indirizzi IPv6}
\label{sec:sock_sa_ipv6}

Essendo IPv6 un'estensione di IPv4, i socket di tipo \const{AF\_INET6} sono
sostanzialmente identici ai precedenti; la parte in cui si trovano
praticamente tutte le differenze fra i due socket è quella della struttura
degli indirizzi; la sua definizione, presa da \headfile{netinet/in.h}, è
riportata in fig.~\ref{fig:sock_sa_ipv6_struct}.

\begin{figure}[!htb]
  \footnotesize \centering
  \begin{minipage}[c]{0.80\textwidth}
    \includestruct{listati/sockaddr_in6.h}
  \end{minipage} 
  \caption{La struttura \structd{sockaddr\_in6} degli indirizzi dei socket
    IPv6 e la struttura \structd{in6\_addr} degli indirizzi IPv6.}
  \label{fig:sock_sa_ipv6_struct}
\end{figure}

Il campo \var{sin6\_family} deve essere sempre impostato ad \constd{AF\_INET6},
il campo \var{sin6\_port} è analogo a quello di IPv4 e segue le stesse regole;
il campo \var{sin6\_flowinfo} è a sua volta diviso in tre parti di cui i 24
bit inferiori indicano l'etichetta di flusso, i successivi 4 bit la priorità e
gli ultimi 4 sono riservati. Questi valori fanno riferimento ad alcuni campi
specifici dell'header dei pacchetti IPv6 (vedi sez.~\ref{sec:IP_ipv6head}) ed
il loro uso è sperimentale.

Il campo \var{sin6\_addr} contiene l'indirizzo a 128 bit usato da IPv6,
espresso da un vettore di 16 byte; anche in questo caso esistono alcuni valori
predediniti, ma essendo il campo un vettore di byte non è possibile assegnarli
con il calore di una costante. Esistono però le variabili predefinite
\var{in6addr\_any} (che indica l'indirizzo generico) e \var{in6addr\_loopback}
(che indica l'indirizzo di loopback) il cui valore può essere copiato in
questo campo. A queste due variabili si aggiungono le macro
\macrod{IN6ADDR\_ANY\_INIT} e \macrod{IN6ADDR\_LOOPBACK\_INIT} per effettuare
delle assegnazioni statiche.

Infine il campo \var{sin6\_scope\_id} è un campo introdotto in Linux con il
kernel 2.4, per gestire alcune operazioni riguardanti il
\textit{multicasting}, è supportato solo per gli indirizzi di tipo
\textit{link-local} (vedi sez.~\ref{sec:IP_ipv6_unicast}) e deve contenere
l'\textit{interface index} (vedi sez.~\ref{sec:sock_ioctl_netdevice}) della
scheda di rete.  Si noti infine che \struct{sockaddr\_in6} ha una dimensione
maggiore della struttura \struct{sockaddr} generica di
fig.~\ref{fig:sock_sa_gen_struct}, quindi occorre stare attenti a non avere
fatto assunzioni riguardo alla possibilità di contenere i dati nelle
dimensioni di quest'ultima (per questo se necessario è opportuno usare
\struct{sockaddr\_storage}).


\subsection{La struttura degli indirizzi locali}
\label{sec:sock_sa_local}

I socket di tipo \const{AF\_UNIX} o \const{AF\_LOCAL} vengono usati per una
comunicazione fra processi che stanno sulla stessa macchina (per questo
vengono chiamati \textit{local domain} o anche \textit{Unix domain}); essi
hanno la caratteristica ulteriore di poter essere creati anche in maniera
anonima attraverso la funzione \func{socketpair} (che abbiamo trattato in
sez.~\ref{sec:ipc_socketpair}).  Quando però si vuole fare riferimento
esplicito ad uno di questi socket si deve usare una struttura degli indirizzi
di tipo \struct{sockaddr\_un}, la cui definizione si è riportata in
fig.~\ref{fig:sock_sa_local_struct}.

\begin{figure}[!htb]
  \footnotesize \centering
  \begin{minipage}[c]{0.80\textwidth}
    \includestruct{listati/sockaddr_un.h}
  \end{minipage} 
  \caption{La struttura \structd{sockaddr\_un} degli indirizzi dei socket
    locali (detti anche \textit{unix domain}) definita in
    \headfiled{sys/un.h}.}
  \label{fig:sock_sa_local_struct}
\end{figure}

In questo caso il campo \var{sun\_family} deve essere \constd{AF\_UNIX},
mentre il campo \var{sun\_path} deve specificare un indirizzo. Questo ha due
forme; può essere ``\textit{named}'' ed in tal caso deve corrispondere ad un
file (di tipo socket) presente nel filesystem o essere ``\textit{abstract}''
nel qual caso viene identificato da una stringa univoca in uno spazio di nomi
astratto. 

Nel primo caso l'indirizzo viene specificato in \var{sun\_path} come una
stringa (terminata da uno zero) corrispondente al \textit{pathname} del file;
nel secondo caso (che è specifico di Linux e non portabile) \var{sun\_path}
deve iniziare con uno zero ed il nome verrà costituito dai restanti byte che
verranno interpretati come stringa senza terminazione (un byte nullo non ha in
questo caso nessun significato).

In realtà esiste una terza forma, \textit{unnamed}, che non è possibile
indicare in fase di scrittura, ma che è quella che viene usata quando si legge
l'indirizzo di un socket anonimo creato con \texttt{socketpair}; in tal caso
la struttura restituita è di dimensione \code{sizeof(sa\_family\_t)}, quindi
\var{sun\_path} non esiste e non deve essere referenziato.

\subsection{La struttura degli indirizzi AppleTalk}
\label{sec:sock_sa_appletalk}

I socket di tipo \const{AF\_APPLETALK} sono usati dalla libreria
\file{netatalk} per implementare la comunicazione secondo il protocollo
AppleTalk, uno dei primi protocolli di rete usato nel mondo dei personal
computer, usato dalla Apple per connettere fra loro computer e stampanti. Il
kernel supporta solo due strati del protocollo, DDP e AARP, e di norma è
opportuno usare le funzioni della libreria \texttt{netatalk}, tratteremo qui
questo argomento principalmente per mostrare l'uso di un protocollo
alternativo.

I socket AppleTalk permettono di usare il protocollo DDP, che è un protocollo
a pacchetto, di tipo \const{SOCK\_DGRAM}; l'argomento \param{protocol} di
\func{socket} deve essere nullo. È altresì possibile usare i socket raw
specificando un tipo \const{SOCK\_RAW}, nel qual caso l'unico valore valido
per \param{protocol} è \constd{ATPROTO\_DDP}.

Gli indirizzi AppleTalk devono essere specificati tramite una struttura
\struct{sockaddr\_atalk}, la cui definizione è riportata in
fig.~\ref{fig:sock_sa_atalk_struct}; la struttura viene dichiarata includendo
il file \headfiled{netatalk/at.h}.

\begin{figure}[!htb]
  \footnotesize \centering
  \begin{minipage}[c]{0.80\textwidth}
    \includestruct{listati/sockaddr_atalk.h}
  \end{minipage} 
  \caption{La struttura \structd{sockaddr\_atalk} degli indirizzi dei socket
    AppleTalk, e la struttura \structd{at\_addr} degli indirizzi AppleTalk.}
  \label{fig:sock_sa_atalk_struct}
\end{figure}

Il campo \var{sat\_family} deve essere sempre \constd{AF\_APPLETALK}, mentre
il campo \var{sat\_port} specifica la porta che identifica i vari
servizi. Valori inferiori a 129 sono usati per le \textsl{porte riservate}, e
possono essere usati solo da processi con i privilegi di amministratore o con
la \textit{capability} \const{CAP\_NET\_BIND\_SERVICE}.  

L'indirizzo remoto è specificato nella struttura \var{sat\_addr}, e deve
essere in \textit{network order} (vedi sez.~\ref{sec:endianness}); esso è
composto da un parte di rete data dal campo \var{s\_net}, che può assumere il
valore \constd{AT\_ANYNET}, che indica una rete generica e vale anche per
indicare la rete su cui si è, il singolo nodo è indicato da \var{s\_node}, e
può prendere il valore generico \constd{AT\_ANYNODE} che indica anche il nodo
corrente, ed il valore \constd{ATADDR\_BCAST} che indica tutti i nodi della
rete.


\subsection{La struttura degli indirizzi dei \textit{packet socket}}
\label{sec:sock_sa_packet}

I \textit{packet socket}, identificati dal dominio \const{AF\_PACKET}, sono
un'interfaccia specifica di Linux per inviare e ricevere pacchetti
direttamente su un'interfaccia di rete, senza passare per le funzioni di
gestione dei protocolli di livello superiore. In questo modo è possibile
implementare dei protocolli in \textit{user space}, agendo direttamente sul
livello fisico. In genere comunque si preferisce usare la libreria
\file{pcap},\footnote{la libreria è mantenuta insieme al comando
  \cmd{tcpdump}, informazioni e documentazione si possono trovare sul sito del
  progetto \url{http://www.tcpdump.org/}.}  che assicura la portabilità su
altre piattaforme, anche se con funzionalità ridotte.

Questi socket possono essere di tipo \const{SOCK\_RAW} o \const{SOCK\_DGRAM}.
Con socket di tipo \const{SOCK\_RAW} si può operare sul livello di
collegamento, ed i pacchetti vengono passati direttamente dal socket al driver
del dispositivo e viceversa.  In questo modo, in fase di trasmissione, il
contenuto completo dei pacchetti, comprese le varie intestazioni, deve essere
fornito dall'utente. In fase di ricezione invece tutto il contenuto del
pacchetto viene passato inalterato sul socket, anche se il kernel analizza
comunque il pacchetto, riempiendo gli opportuni campi della struttura
\struct{sockaddr\_ll} ad esso associata.

Si usano invece socket di tipo \const{SOCK\_DGRAM} quando si vuole operare a
livello di rete. 

In questo caso in fase di ricezione l'intestazione del protocollo di
collegamento viene rimossa prima di passare il resto del pacchetto all'utente,
mentre in fase di trasmissione viene creata una opportuna intestazione per il
protocollo a livello di collegamento utilizzato, usando le informazioni
necessarie che devono essere specificate sempre con una struttura
\struct{sockaddr\_ll}.

Nella creazione di un \textit{packet socket} il valore dell'argomento
\param{protocol} di \func{socket} serve a specificare, in \textit{network
  order}, il numero identificativo del protocollo di collegamento si vuole
utilizzare. I valori possibili sono definiti secondo lo standard IEEE 802.3, e
quelli disponibili in Linux sono accessibili attraverso opportune costanti
simboliche definite nel file \file{linux/if\_ether.h}. Se si usa il valore
speciale \constd{ETH\_P\_ALL} passeranno sul \textit{packet socket} tutti i
pacchetti, qualunque sia il loro protocollo di collegamento. Ovviamente l'uso
di questi socket è una operazione privilegiata e può essere effettuati solo da
un processo con i privilegi di amministratore (\ids{UID} effettivo nullo) o
con la \textit{capability} \const{CAP\_NET\_RAW}.

Una volta aperto un \textit{packet socket}, tutti i pacchetti del protocollo
specificato passeranno attraverso di esso, qualunque sia l'interfaccia da cui
provengono; se si vuole limitare il passaggio ad una interfaccia specifica
occorre usare la funzione \func{bind} (vedi sez.~\ref{sec:TCP_func_bind}) per
agganciare il socket a quest'ultima.

\begin{figure}[!htb]
  \footnotesize \centering
  \begin{minipage}[c]{\textwidth}
    \includestruct{listati/sockaddr_ll.h}
  \end{minipage} 
  \caption{La struttura \structd{sockaddr\_ll} degli indirizzi dei
    \textit{packet socket}.}
  \label{fig:sock_sa_packet_struct}
\end{figure}

Nel caso dei \textit{packet socket} la struttura degli indirizzi è di tipo
\struct{sockaddr\_ll}, e la sua definizione è riportata in
fig.~\ref{fig:sock_sa_packet_struct}; essa però viene ad assumere un ruolo
leggermente diverso rispetto a quanto visto finora per gli altri tipi di
socket.  Infatti se il socket è di tipo \const{SOCK\_RAW} si deve comunque
scrivere tutto direttamente nel pacchetto, quindi la struttura non serve più a
specificare gli indirizzi. Essa mantiene questo ruolo solo per i socket di
tipo \const{SOCK\_DGRAM}, per i quali permette di specificare i dati necessari
al protocollo di collegamento, mentre viene sempre utilizzata in lettura (per
entrambi i tipi di socket), per la ricezione dei dati relativi a ciascun
pacchetto.

Al solito il campo \var{sll\_family} deve essere sempre impostato al valore
\constd{AF\_PACKET}. Il campo \var{sll\_protocol} indica il protocollo scelto,
e deve essere indicato in \textit{network order}, facendo uso delle costanti
simboliche definite in \file{linux/if\_ether.h}. Il campo \var{sll\_ifindex} è
l'indice dell'interfaccia (l'\textit{inxterface index} (vedi
sez.~\ref{sec:sock_ioctl_netdevice}) che in caso di presenza di più
interfacce dello stesso tipo (se ad esempio si hanno più schede Ethernet),
permette di selezionare quella con cui si vuole operare (un valore nullo
indica qualunque interfaccia).  Questi sono i due soli campi che devono essere
specificati quando si vuole selezionare una interfaccia specifica, usando
questa struttura con la funzione \func{bind}.

I campi \var{sll\_halen} e \var{sll\_addr} indicano rispettivamente
l'indirizzo associato all'interfaccia sul protocollo di collegamento e la
relativa lunghezza; ovviamente questi valori cambiano a seconda del tipo di
collegamento che si usa, ad esempio, nel caso di Ethernet, questi saranno il
MAC address della scheda e la relativa lunghezza. Essi vengono usati, insieme
ai campi \var{sll\_family} e \var{sll\_ifindex} quando si inviano dei
pacchetti, in questo caso tutti gli altri campi devono essere nulli.

Il campo \var{sll\_hatype} indica il tipo ARP, come definito in
\file{linux/if\_arp.h}, mentre il campo \var{sll\_pkttype} indica il tipo di
pacchetto; entrambi vengono impostati alla ricezione di un pacchetto ed han
senso solo in questo caso. In particolare \var{sll\_pkttype} può assumere i
seguenti valori: \constd{PACKET\_HOST} per un pacchetto indirizzato alla
macchina ricevente, \constd{PACKET\_BROADCAST} per un pacchetto di
\textit{broadcast}, \constd{PACKET\_MULTICAST} per un pacchetto inviato ad un
indirizzo fisico di \textit{multicast}, \constd{PACKET\_OTHERHOST} per un
pacchetto inviato ad un'altra stazione (e ricevuto su un'interfaccia in modo
promiscuo), \constd{PACKET\_OUTGOING} per un pacchetto originato dalla propria
macchina che torna indietro sul socket.

Si tenga presente infine che in fase di ricezione, anche se si richiede il
troncamento del pacchetto, le funzioni \func{recv}, \func{recvfrom} e
\func{recvmsg} (vedi sez.~\ref{sec:net_sendmsg}) restituiranno comunque la
lunghezza effettiva del pacchetto così come arrivato sulla linea.

%% \subsection{La struttura degli indirizzi DECnet}
%% \label{sec:sock_sa_decnet}
 
%% I socket di tipo \const{AF\_DECnet} usano il protocollo DECnet, usato dai VAX
%% Digital sotto VMS quando ancora il TCP/IP non era diventato lo standard di
%% fatto. Il protocollo è un protocollo chiuso, ed il suo uso attuale è limitato
%% alla comunicazione con macchine che stanno comunque scomparendo. Lo si riporta
%% solo come esempio 


% TODO aggiungere AF_CAN, vedi http://lwn.net/Articles/253425, dal 2.6.25 ?

% TODO: trattare i socket RDS, vedi documentazione del kernel, file 
% Documentation/networking/rds.txt



\section{Le funzioni di conversione degli indirizzi}
\label{sec:sock_addr_func}

In questa sezione tratteremo delle varie funzioni usate per manipolare gli
indirizzi, limitandoci però agli indirizzi Internet.  Come accennato gli
indirizzi e i numeri di porta usati nella rete devono essere forniti nel
cosiddetto \textit{network order}, che corrisponde al formato \textit{big
  endian} (vedi sez.~\ref{sec:endianness}), anche quando la proprio macchina
non usa questo formato, cosa che può comportare la necessità di eseguire delle
conversioni.


\subsection{Le funzioni per il riordinamento}
\label{sec:sock_func_ord}

Come già visto in sez.~\ref{sec:endianness} il problema connesso
all'\textit{endianness} è che quando si passano dei dati da un tipo di
architettura all'altra i dati vengono interpretati in maniera diversa, e ad
esempio nel caso dell'intero a 16 bit ci si ritroverà con i due byte in cui è
suddiviso scambiati di posto.  

Per questo motivo si usano delle funzioni di conversione che servono a tener
conto automaticamente della possibile differenza fra l'ordinamento usato sul
computer e quello che viene usato nelle trasmissione sulla rete; queste
funzioni sono \funcd{htonl}, \funcd{htons}, \funcd{ntohl} e \funcd{ntohs} ed i
rispettivi prototipi sono:

\begin{funcproto}{
\fhead{arpa/inet.h}
\fdecl{unsigned long int htonl(unsigned long int hostlong)}
\fdesc{Converte l'intero a 32 bit \param{hostlong} dal formato della macchina a
  quello della rete.} 
\fdecl{unsigned short int htons(unsigned short int hostshort)}
\fdesc{Converte l'intero a 16 bit \param{hostshort} dal formato della macchina a
  quello della rete.}
\fdecl{unsigned long int ntohl(unsigned long int netlong)}
\fdesc{Converte l'intero a 32 bit \param{netlong} dal formato della rete a
  quello della macchina.}
\fdecl{unsigned sort int ntohs(unsigned short int netshort)}
\fdesc{Converte l'intero a 16 bit \param{netshort} dal formato della rete a
  quello della macchina.}
}

{Tutte le funzioni restituiscono il valore convertito, e non prevedono
  errori.}
\end{funcproto}

I nomi sono assegnati usando la lettera \texttt{n} come mnemonico per indicare
l'ordinamento usato sulla rete (da \textit{network order}) e la lettera
\texttt{h} come mnemonico per l'ordinamento usato sulla macchina locale (da
\textit{host order}), mentre le lettere \texttt{s} e \texttt{l} stanno ad
indicare i tipi di dato (\ctyp{long} o \ctyp{short}, riportati anche dai
prototipi).

Usando queste funzioni si ha la conversione automatica: nel caso in cui la
macchina che si sta usando abbia una architettura \textit{big endian} queste
funzioni sono definite come macro che non fanno nulla. Per questo motivo vanno
sempre utilizzate, anche quando potrebbero non essere necessarie, in modo da
assicurare la portabilità del codice su tutte le architetture.


\subsection{Le funzioni di conversione per gli indirizzi IPv4}
\label{sec:sock_func_ipv4}

Un secondo insieme di funzioni di manipolazione è quello che serve per passare
dalla rappresentazione simbolica degli indirizzi IP al formato binario
previsto dalla struttura degli indirizzi di
fig.~\ref{fig:sock_sa_ipv4_struct}, e viceversa. La notazione più comune è la
cosiddetta notazione \itindex{dotted-decimal} \textit{dotted-decimal}, che
prevede che gli indirizzi IPv4 siano indicati con l'espressione del valore
numerico decimale di ciascuno dei 4 byte che li costituiscono separati da un
punto (ad esempio \texttt{192.168.0.1}).

In realtà le funzioni che illustreremo supportano una notazione che più
propriamente dovrebbe esser chiamata \textit{numbers-and-dot} in quanto il
valore può essere indicato con numeri espressi sia in decimale, che in ottale
(se indicati apponendo uno zero) che in esadecimale (se indicati apponendo
\texttt{0x}). Inoltre per la parte meno significativa dell'espressione, quella
che riguarda l'indirizzo locale, si può usare, eliminando altrettanti punti,
valori a 16 o a 24 bit, e togliendo tutti i punti, si può usare anche
direttamente un valore numerico a 32 bit.\footnote{la funzionalità si trova
  anche in gran parte dei programmi che usano indirizzi di rete, e deriva
  direttamente da queste funzioni.}

Tradizionalmente la conversione di un indirizzo \textit{dotted-decimal} al
valore numerico veniva eseguita dalla funzione \funcd{inet\_addr} (prevista
fin dalle origini in BSD e inclusa in POSIX.1-2001) il cui prototipo è:

\begin{funcproto}{
\fhead{arpa/inet.h}
\fdecl{in\_addr\_t inet\_addr(const char *strptr)}
\fdesc{Converte la stringa dell'indirizzo \textit{dotted decimal} in nel
  numero IP in network order.} 
}

{La funzione ritorna il valore dell'indirizzo in caso di successo e
  \const{INADDR\_NONE} per un errore e non genera codici di errore.}
\end{funcproto}

La prima funzione, \func{inet\_addr}, restituisce l'indirizzo a 32 bit in
\textit{network order} (del tipo \type{in\_addr\_t}) a partire dalla stringa
passata nell'argomento \param{strptr}. In caso di errore (quando la stringa
non esprime un indirizzo valido) restituisce invece il valore
\const{INADDR\_NONE}, che tipicamente sono trentadue bit a uno.  Questo però
comporta che la stringa \texttt{255.255.255.255}, che pure è un indirizzo
valido, non può essere usata con questa funzione dato che genererebe comunque
un errore; per questo motivo essa è generalmente deprecata in favore di
\func{inet\_aton}.

Per effettuare la conversione inversa la funzione usata tradizionalmente è
\funcd{inet\_ntoa}, anch'essa presente fin da BSD 4.3, in cui si riprende la
notazione già vista in sez.~\ref{sec:sock_func_ord} che usa la lettera
\texttt{n} come mnemonico per indicare la rete ed \texttt{a} (per ASCII) come
mnemonico per indicare la stringa corrispodente all'indirizzo; il suo
prototipo è:

\begin{funcproto}{
\fhead{arpa/inet.h}
\fdecl{char *inet\_ntoa(struct in\_addr addrptr)}
\fdesc{Converte un indirizzo IP in una stringa \textit{dotted decimal}.} 
}

{La funzione l'indirizzo della stringa con il valore dell'indirizzo convertito
  e non prevede errori.}
\end{funcproto}

La funzione converte il valore a 32 bit dell'indirizzo, espresso in
\textit{network order}, e preso direttamente con un puntatore al relativo
campo della struttura degli indirizzi, restituendo il puntatore alla stringa
che contiene l'espressione in formato \textit{dotted-decimal}. Si deve tenere
presente che la stringa risiede in un segmento di memoria statica, per cui
viene riscritta ad ogni chiamata e la funzione non è rientrante.

Per rimediare ai problemi di \funcd{inet\_addr} è stata sostituita da
\funcd{inet\_aton}, che però non è stata standardizzata e non è presente in
POSIX.1-2001, anche se è definita sulla gran parte dei sistemi Unix; il suo
prototipo è:

\begin{funcproto}{
\fhead{arpa/inet.h}
\fdecl{int inet\_aton(const char *src, struct in\_addr *dest)}
\fdesc{Converte la stringa dell'indirizzo \textit{dotted decimal} in un
  indirizzo IP.}
}

{La funzione ritorna un valore non nullo in caso di successo e $0$ per un
  errore e non genera codici di errore.}
\end{funcproto}

La funzione converte la stringa puntata da \param{src} nell'indirizzo binario
che viene memorizzato nell'opportuna struttura \struct{in\_addr} (si veda
fig.~\ref{fig:sock_sa_ipv4_struct}) situata all'indirizzo dato
dall'argomento \param{dest} (è espressa in questa forma in modo da poterla
usare direttamente con il puntatore usato per passare la struttura degli
indirizzi). La funzione restituisce un valore diverso da zero se l'indirizzo è
valido e la conversione ha successo e 0 in caso contrario. Se usata
con \param{dest} inizializzato a \val{NULL} può essere usata per effettuare la
validazione dell'indirizzo espresso da \param{src}.

Oltre a queste tre funzioni esistono le ulteriori \funcm{inet\_lnaof},
\funcm{inet\_netof} e \funcm{inet\_makeaddr} che assumono la ormai obsoleta e
deprecata suddivisione in classi degli indirizzi IP per fornire la parte di
rete e quella di indirizzo locale. Ad oggi il loro uso non ha più alcun senso
per ciò non le tratteremo.


\subsection{Le funzioni di conversione per indirizzi IP generici}
\label{sec:sock_conv_func_gen}

Le tre funzioni precedenti sono limitate solo ad indirizzi IPv4, per questo
motivo è preferibile usare le due nuove funzioni \func{inet\_pton} e
\func{inet\_ntop} che possono convertire anche gli indirizzi IPv6. Anche in
questo caso le lettere \texttt{n} e \texttt{p} sono degli mnemonici per
ricordare il tipo di conversione effettuata e stanno per \textit{presentation}
e \textit{numeric}.

Entrambe le funzioni accettano l'argomento \param{af} che indica il tipo di
indirizzo, e che può essere soltanto \const{AF\_INET} o \const{AF\_INET6}. La
prima funzione, \funcd{inet\_pton}, serve a convertire una stringa in un
indirizzo; il suo prototipo è:

\begin{funcproto}{
\fhead{sys/socket.h}
\fdecl{int inet\_pton(int af, const char *src, void *addr\_ptr)} 
\fdesc{Converte l'indirizzo espresso tramite una stringa nel valore numerico.} 
}

{La funzione ritorna $1$ in caso di successo, $0$ se \param{src} non contiene
  una rappresentazione valida per la famiglia di indirizzi indicati
  da \param{af} e $-1$ se \param{af} specifica una famiglia di indirizzi non
  valida, e solo in quest'ultimo caso  \var{errno} assumerà il valore
  \errcode{EAFNOSUPPORT}.
}
\end{funcproto}

La funzione converte la stringa indicata tramite \param{src} nel valore
numerico dell'indirizzo IP del tipo specificato da \param{af} che viene
memorizzato all'indirizzo puntato da \param{addr\_ptr}. La funzione supporta
per IPv4 la sola notazione \textit{dotted-decimal}, e non quella più completa
\textit{number-and-dot} che abbiamo visto per \func{inet\_aton}. Per IPv6 la
notazione prevede la suddivisione dei 128 bit dell'indirizzo in 16 parti di 16
bit espresse con valori esadecimali separati dal carattere ``\texttt{:}'' ed
una serie di valori nulli possono essere sostituiti (una sola volta, sempre a
partire dalla sinistra) con la notazione ``\texttt{::}'', un esempio di
indirizzo in questa forma potrebbe essere \texttt{2001:db8::8:ba98:2078:e3e3},
per una descrizione più completa si veda sez.~\ref{sec:IP_ipv6_notation}.

La seconda funzione di conversione è \funcd{inet\_ntop} che converte un
indirizzo in una stringa; il suo prototipo è:

\begin{funcproto}{
\fhead{sys/socket.h}
\fdecl{char *inet\_ntop(int af, const void *addr\_ptr, char *dest, size\_t len)}
\fdesc{Converte l'indirizzo dalla relativa struttura in una stringa simbolica.} 
}

{La funzione ritorna un puntatore non nullo alla stringa convertita in caso di
  successo e \val{NULL} per un errore, nel qual caso \var{errno} assumerà uno
  dei valori:
  \begin{errlist}
    \item[\errcode{ENOSPC}] le dimensioni della stringa con la conversione
      dell'indirizzo eccedono la lunghezza specificata da \param{len}.
    \item[\errcode{ENOAFSUPPORT}] la famiglia di indirizzi \param{af} non è
      una valida.
  \end{errlist}
}
\end{funcproto}


La funzione converte la struttura dell'indirizzo puntata da \param{addr\_ptr}
in una stringa che viene copiata nel buffer puntato dall'indirizzo
\param{dest}; questo deve essere preallocato dall'utente e la lunghezza deve
essere almeno \constd{INET\_ADDRSTRLEN} in caso di indirizzi IPv4 e
\constd{INET6\_ADDRSTRLEN} per indirizzi IPv6; la lunghezza del buffer deve
comunque essere specificata con il parametro \param{len}.

Gli indirizzi vengono convertiti da/alle rispettive strutture di indirizzo
(una struttura \struct{in\_addr} per IPv4, e una struttura \struct{in6\_addr}
per IPv6), che devono essere precedentemente allocate e passate attraverso il
puntatore \param{addr\_ptr}; l'argomento \param{dest} di \func{inet\_ntop} non
può essere nullo e deve essere allocato precedentemente.





% LocalWords:  socket sez cap BSD SVr XTI Transport Interface TCP stream UDP PF
% LocalWords:  datagram broadcast descriptor sys int domain type protocol errno
% LocalWords:  EPROTONOSUPPORT ENFILE kernel EMFILE EACCES EINVAL ENOBUFS raw
% LocalWords:  ENOMEM table family AF address name glibc UNSPEC LOCAL Local IPv
% LocalWords:  communication INET protocols ip AX Amateur IPX Novell APPLETALK
% LocalWords:  Appletalk ddp NETROM NetROM Multiprotocol ATMPVC Access to ATM
% LocalWords:  PVCs ITU ipv PLP DECnet Reserved for project NETBEUI LLC KEY key
% LocalWords:  SECURITY Security callback NETLINK interface device netlink Low
% LocalWords:  PACKET level packet ASH Ash ECONET Acorn Econet ATMSVC SVCs SNA
% LocalWords:  IRDA PPPOX PPPoX WANPIPE Wanpipe BLUETOOTH Bluetooth POSIX bits
% LocalWords:  dall'header tab SOCK capabilities capability styles DGRAM read
% LocalWords:  SEQPACKET RDM sockaddr reference void fig Header uint socklen at
% LocalWords:  addr netinet port len Stevens unsigned short casting nell'header
% LocalWords:  BIND SERVICE bind union order big endian flowinfo dell'header ll
% LocalWords:  multicast multicasting local socketpair sun path filesystem AARP
% LocalWords:  pathname AppleTalk netatalk personal Apple ATPROTO atalk sat if
% LocalWords:  ANYNET node ANYNODE ATADDR BCAST pcap IEEE linux ether ETH ALL
% LocalWords:  sll ifindex ethernet halen MAC hatype ARP arp pkttype HOST recv
% LocalWords:  OTHERHOST OUTGOING recvfrom recvmsg endianness little endtest Mac
% LocalWords:  Intel Digital Motorola IBM VME PowerPC l'Intel xABCD ptr htonl
% LocalWords:  htons ntohl ntohs long hostlong hostshort netlong
% LocalWords:  sort netshort host inet aton ntoa dotted decimal const char src
% LocalWords:  strptr struct dest addrptr INADDR NULL pton ntop presentation af
% LocalWords:  numeric EAFNOSUPPORT size ENOSPC ENOAFSUPPORT ADDRSTRLEN ROUTE
% LocalWords:  of tcpdump page


%%% Local Variables: 
%%% mode: latex
%%% TeX-master: "gapil"
%%% End: 
